\section{Task and Basic Concept}

\noindent This project presents an autonomous "Block Assembler Robot" with the task of building a structure of magnetic blocks on top of a 6x6 grid of "base blocks". The structure is not hard coded, but rather prespecified by a JSON file (an example is given in Appendix A -- Section \ref{sec:appendix}), which contains the grid coordinates and colors of every block in the structure. The robot is able to go from almost any existing structure state to almost any JSON-specified structure state. This invariance of initial state means that the system is robust to most user actions including users adding and removing blocks from the structure or connecting and moving blocks around in the pile because the system is memory-less and therefore determines the next action to get the structure state closer to the desired final state every iteration of adding/removing a block.

\noindent Due to the open-ended nature of the structure the robot will be asked to build, there are a number of challenges risen by the task. The primary difficulties include accurately determining the state of the structure at any state despite the fact that some blocks may be obscured by others, as well as robustly recovering from missteps (finding misplaced blocks and placing them where they are intended to go) without further damaging the structure. 
